%% =====================================================================
%%  Talleres GARyC -- Análisis Económico y Computacional
%%  Syllabus
%%  Compilar con pdfLaTeX (Overleaf: Menú > Compilador > pdfLaTeX).
%% =====================================================================
\documentclass[11pt]{article}
\usepackage{garyc}

\documento{Syllabus}
\title{Talleres GARyC: Análisis Económico y Computacional. Syllabus}
\author{Talleres GARyC}
\date{}

% Abreviaturas de uso frecuente en la bibliografía
\newcommand{\PUP}{Princeton University Press}
\newcommand{\CUP}{Cambridge University Press}
\newcommand{\CRC}{Chapman \& Hall/CRC}
\newcommand{\JoE}{\textit{Journal of Econometrics}}
\newcommand{\AngristMIT}[1]{\obra{Angrist}{2021}{Econometric Data Science, notas de clase 14.320}{MIT. #1}}

\begin{document}

\portada{Syllabus}

%% ---------------------------------------------------------------------
\begin{presentacion}
Los talleres tienen un formato principalmente práctico: una breve
introducción conceptual y el grueso de la sesión dedicado a implementar
cada método en vivo sobre datos o problemas reales de la gerencia. El
hilo conductor es una sola idea:
\[
  \text{diferencia observada} \;=\; \text{efecto causal} \;+\; \text{sesgo de selección}.
\]
Cada diseño de evaluación es una manera distinta de enfrentar ese sesgo.

Las herramientas de IA generativa se incorporan como eje transversal en
todos los módulos, integradas al flujo de trabajo: explicación de
conceptos, generación y depuración de código, interpretación de
resultados y apoyo en la documentación. Cada módulo incluye su
bibliografía de referencia, basada en las respuestas a la encuesta
diagnóstica y en las recomendaciones del equipo.

El contenido se organiza en tres talleres: los dos primeros de 20 horas
cada uno y el tercero de 13 horas, a razón de una hora por sesión. El
Taller~1 reúne lo sustantivo para evaluar el impacto de una política;
los Talleres~2 y~3 profundizan y complementan. Cada tema se presenta en
tres capas: intuición para todos, formalización opcional y práctica con
código, apoyadas en datos (sintéticos) de temática bancaria y
aplicaciones equivalentes en R, Python y Stata.
\end{presentacion}

%% ---------------------------------------------------------------------
\clearpage
\section*{Estructura del programa}
\addcontentsline{toc}{section}{Estructura del programa}

\begin{table}[ht]
\centering
\small
\renewcommand{\arraystretch}{1.25}
\caption{Talleres, carga horaria y módulos.}
\label{tab:estructura}
\begin{tabularx}{\linewidth}{@{}P{1.55cm} P{1.0cm} P{3.6cm} L@{}}
\toprule
\sffamily\bfseries Taller & \sffamily\bfseries Horas & \sffamily\bfseries Enfoque & \sffamily\bfseries Módulos (horas) \\
\midrule
\textbf{Taller 1} & 20 h &
Evaluación de impacto: fundamentos y diseños principales (núcleo) &
M1~Repaso exprés de estadística inferencial (1)\newline
M2~Marco causal y experimentos aleatorizados (5)\newline
M3~Regresión como descripción de la función de esperanza condicional (CEF) (3)\newline
M4~Regresión causal: OVB y CIA (2)\newline
M5~Inferencia aplicada (1)\newline
M6~Datos panel y diferencias en diferencias (4)\newline
M7~Variables instrumentales y regresión discontinua (2)\newline
M8~Síntesis y taller de proyectos (2) \\
\midrule
\textbf{Taller 2} & 20 h &
Econometría aplicada y diseños complementarios &
M1~Modelos de elección y variables limitadas (3)\newline
M2~Selección, matching y sensibilidad (3)\newline
M3~Econometría de panel avanzada (2)\newline
M4~Diferencias en diferencias avanzado y control sintético (4)\newline
M5~Experimentos avanzados (3)\newline
M6~Variables instrumentales y regresión discontinua avanzados (3)\newline
M7~Aprendizaje automático causal (2) \\
\midrule
\textbf{Taller 3} & 13 h &
Análisis geoespacial y actividad de aplicación &
M1~Análisis geoespacial (7)\newline
M2~Evaluación con dimensión espacial (3)\newline
M3~Actividad con aplicación a temas de la gerencia (3) \\
\bottomrule
\end{tabularx}
\end{table}

%% =====================================================================
\taller{1}{Evaluación de impacto: fundamentos y diseños principales}
  {20 horas \textperiodcentered{} 20 sesiones de aproximadamente 1 hora \textperiodcentered{} núcleo del programa}

%% ---------------------------------------------------------------------
\modulo{1}{Repaso exprés de estadística inferencial}{1 h}
\begin{sesiones}
\sesion{1} De la muestra a la inferencia: pruebas de hipótesis,
  comparación de medias y proporciones, correlación y regresión lineal simple.
\end{sesiones}
\begin{bibliografia}
\obra{Casella y Berger}{2002}{Statistical Inference}{2.ª ed., Duxbury.} \recomendado
\obra{Kleiber y Zeileis}{2008}{Applied Econometrics with R}{Springer.}
\end{bibliografia}

%% ---------------------------------------------------------------------
\modulo{2}{Marco causal y experimentos aleatorizados}{5 h}
\begin{sesiones}
\sesion{2} De la pregunta de política al estimando: unidad, tratamiento,
  resultado, horizonte y contrafactual; ATE, ATT, ITT y LATE; validez interna
  y externa. Lienzo de diseño de evaluación.
\sesion{3} Resultados potenciales y sesgo de selección: diferencia observada
  = efecto sobre los tratados + sesgo de selección; por qué la aleatorización
  lo elimina. Simulación con $Y_{1}$ y $Y_{0}$ conocidos.
\sesion{4} Grafos causales o grafos acíclicos dirigidos (DAGs): confusores,
  mediadores y colisionadores; qué controlar y qué no.
\sesion{5} Experimentos aleatorizados I: tablas de balance, estimación con
  diferencia de medias y con regresión, aleatorización estratificada. Caso:
  estado de cuenta estandarizado de tarjetas de crédito (datos sintéticos o
  públicos).
\sesion{6} Experimentos aleatorizados II: poder y efecto mínimo detectable,
  cumplimiento imperfecto (ITT), atrición, efectos indirectos, ética y
  pre-registro.
\end{sesiones}
\begin{bibliografia}
\obra{Angrist y Pischke}{2015}{Mastering 'Metrics}{\PUP.}
\obra{Angrist y Pischke}{2009}{Mostly Harmless Econometrics}{\PUP.}
\obra{Huntington-Klein}{2022}{The Effect}{\CRC.}
\obra{Cunningham}{2021}{Causal Inference: The Mixtape}{Yale University Press.}
\obra{List}{próx.}{Experimental Economics: Theory and Practice}{Aplicado a experimentos aleatorizados.}
\AngristMIT{Lecture Note 4.}
\end{bibliografia}

%% ---------------------------------------------------------------------
\modulo{3}{Regresión como descripción de la función de esperanza condicional (CEF)}{3 h}
\begin{sesiones}
\sesion{7} La CEF y la regresión: ley de esperanzas iteradas; regresión como
  mejor aproximación lineal a la CEF; regresión lineal simple y múltiple;
  condicionar = comparar manzanas con manzanas.
\sesion{8} Modelos saturados e interacciones: binarias como diferencias de
  medias; efectos principales e interacciones; residuos, valores ajustados
  y~$R^{2}$.
\sesion{9} Anatomía de la regresión (Frisch--Waugh--Lovell): \enquote{controlar
  por} como limpiar la variación; verificación numérica con regresiones
  auxiliares.
\end{sesiones}
\begin{bibliografia}
\obra{Greene}{2018}{Econometric Analysis}{8.ª ed., Pearson.} \recomendado
\obra{Cameron y Trivedi}{2005}{Microeconometrics: Methods and Applications}{\CUP.}
\obra{Wooldridge}{2010}{Econometric Analysis of Cross Section and Panel Data}{2.ª ed., MIT Press.}
\obra{Kleiber y Zeileis}{2008}{Applied Econometrics with R}{Springer.}
\AngristMIT{Lecture Notes 1, 5, 8 y 9.}
\end{bibliografia}

%% ---------------------------------------------------------------------
\modulo{4}{Regresión causal: OVB y CIA}{2 h}
\begin{sesiones}
\sesion{10} Sesgo por variable omitida (OVB): corta = larga + (efecto de la
  omitida) $\times$ (regresión de la omitida sobre la incluida); signo del
  sesgo y cuándo corta = larga.
\sesion{11} Supuesto de independencia condicional (CIA) y controles buenos y
  malos: lectura de tablas con controles crecientes; variables posteriores al
  tratamiento; primera mirada a la sensibilidad ante variables omitidas.
\end{sesiones}
\begin{bibliografia}
\obra{Angrist y Pischke}{2009}{Mostly Harmless Econometrics}{\PUP.}
\AngristMIT{Lecture Note 6.}
\end{bibliografia}

%% ---------------------------------------------------------------------
\modulo{5}{Inferencia aplicada}{1 h}
\begin{sesiones}
\sesion{12} Errores estándar que sí importan: errores robustos, dependencia
  entre observaciones y clustering (factor de Moulton), pocos clusters.
\end{sesiones}
\begin{bibliografia}
\obra{Angrist y Pischke}{2009}{Mostly Harmless Econometrics}{\PUP{} (cap.~8).}
\obra{Abadie, Athey, Imbens y Wooldridge}{2023}{When Should You Adjust Standard Errors for Clustering?}{\textit{Quarterly Journal of Economics}.}
\AngristMIT{Lecture Notes 7 y 10.}
\end{bibliografia}

%% ---------------------------------------------------------------------
\modulo{6}{Datos panel y diferencias en diferencias}{4 h}
\begin{sesiones}
\sesion{13} Estructura de datos panel, pooled OLS y efectos fijos
  (variación within e interpretación vía FWL).
\sesion{14} Diferencias en diferencias clásico ($2\times2$): tendencias
  paralelas en resultados potenciales, forma de regresión y sensibilidad a la
  escala. Caso: liquidez bancaria en la Gran Depresión.
\sesion{15} Por qué falla el TWFE con tratamiento escalonado: efectos
  heterogéneos en el tiempo y descomposición de Goodman-Bacon.
\sesion{16} Diferencias en diferencias moderno y estudios de evento:
  estimadores robustos (Callaway--Sant'Anna, Sun--Abraham, de
  Chaisemartin--D'Haultf{\oe}uille, imputación); lectura de pre-tendencias.
\end{sesiones}
\begin{bibliografia}
\obra{Wooldridge}{2010}{Econometric Analysis of Cross Section and Panel Data}{2.ª ed., MIT Press.}
\obra{Cameron y Trivedi}{2005}{Microeconometrics: Methods and Applications}{\CUP.}
\obra{Greene}{2018}{Econometric Analysis}{8.ª ed., Pearson.} \recomendado
\obra{Cunningham}{2021}{Causal Inference: The Mixtape}{Yale University Press.}
\obra{Goodman-Bacon}{2021}{Difference-in-Differences with Variation in Treatment Timing}{\JoE.}
\obra{Callaway y Sant'Anna}{2021}{Difference-in-Differences with Multiple Time Periods}{\JoE.}
\obra{Sun y Abraham}{2021}{Estimating Dynamic Treatment Effects in Event Studies with Heterogeneous Treatment Effects}{\JoE.}
\obra{de Chaisemartin y D'Haultf{\oe}uille}{2020}{Two-Way Fixed Effects Estimators with Heterogeneous Treatment Effects}{\textit{American Economic Review}.}
\obra{Roth, Sant'Anna, Bilinski y Poe}{2023}{What's Trending in Difference-in-Differences?}{\JoE.}
\obra{Richardson y Troost}{2009}{Monetary Intervention Mitigated Banking Panics during the Great Depression}{\textit{Journal of Political Economy}.}
\AngristMIT{Lecture Notes 12 y 13.}
\end{bibliografia}

%% ---------------------------------------------------------------------
\modulo{7}{Variables instrumentales y regresión discontinua}{2 h}
\begin{sesiones}
\sesion{17} Variables instrumentales y LATE: primera etapa, forma reducida y
  razón de Wald; complacientes (\textit{compliers}); diseños de aliento.
\sesion{18} Regresión discontinua: RD nítida; umbrales de scoring, tamaño de
  institución o regulatorios; pruebas de balance y de manipulación (primera
  mirada).
\end{sesiones}
\begin{bibliografia}
\obra{Angrist y Pischke}{2009}{Mostly Harmless Econometrics}{\PUP.}
\obra{Angrist y Pischke}{2015}{Mastering 'Metrics}{\PUP.}
\obra{Huntington-Klein}{2022}{The Effect}{\CRC.}
\obra{Cunningham}{2021}{Causal Inference: The Mixtape}{Yale University Press.}
\obra{Imbens y Angrist}{1994}{Identification and Estimation of Local Average Treatment Effects}{\textit{Econometrica}.}
\obra{Imbens y Lemieux}{2008}{Regression Discontinuity Designs: A Guide to Practice}{\JoE.}
\AngristMIT{Lecture Notes 11 y 14.}
\end{bibliografia}

%% ---------------------------------------------------------------------
\modulo{8}{Síntesis y taller de proyectos}{2 h}
\begin{sesiones}
\sesion{19} ¿Qué diseño uso? Árbol de decisión (qué variación tengo
  $\rightarrow$ qué diseño); lectura crítica de una evaluación;
  reproducibilidad y reporte para tomadores de decisión.
\sesion{20} Taller de proyectos: presentación y retroalimentación de los
  lienzos de diseño de la sesión~2, ya refinados.
\end{sesiones}
\begin{bibliografia}
\obra{Huntington-Klein}{2022}{The Effect}{\CRC.}
\obra{Cunningham}{2021}{Causal Inference: The Mixtape}{Yale University Press.}
\obra{List}{2022}{The Voltage Effect}{Currency.}
\end{bibliografia}

%% =====================================================================
\taller{2}{Econometría aplicada y diseños complementarios}
  {20 horas \textperiodcentered{} 20 sesiones de aproximadamente 1 hora \textperiodcentered{} requiere haber tomado el Taller~1}

%% ---------------------------------------------------------------------
\modulo{1}{Modelos de elección y variables limitadas}{3 h}
\begin{sesiones}
\sesion{1} Modelo de probabilidad lineal: ventajas, límites e interpretación
  como aproximación a la CEF de una binaria.
\sesion{2} Modelos Probit y Logit: estimación, efectos marginales y cuándo
  importa elegir entre ellos.
\sesion{3} Modelo Tobit: resultados censurados e interpretación de sus efectos.
\end{sesiones}
\begin{bibliografia}
\obra{Greene}{2018}{Econometric Analysis}{8.ª ed., Pearson.} \recomendado
\obra{Maddala}{1983}{Limited-Dependent and Qualitative Variables in Econometrics}{\CUP.} \recomendado
\obra{Cameron y Trivedi}{2005}{Microeconometrics: Methods and Applications}{\CUP.}
\obra{Wooldridge}{2010}{Econometric Analysis of Cross Section and Panel Data}{2.ª ed., MIT Press.}
\obra{Kleiber y Zeileis}{2008}{Applied Econometrics with R}{Springer.}
\end{bibliografia}

%% ---------------------------------------------------------------------
\modulo{2}{Selección, matching y sensibilidad}{3 h}
\begin{sesiones}
\sesion{4} Matching y ponderación por \textit{propensity score}; relación con
  la regresión y con la CIA.
\sesion{5} Selección sobre no observables: modelo de selección de Heckman;
  restricción de exclusión y cuándo es creíble.
\sesion{6} Sensibilidad ante variables omitidas: estabilidad de coeficientes
  (Oster) y valores de robustez (Cinelli--Hazlett).
\end{sesiones}
\begin{bibliografia}
\obra{Imbens y Rubin}{2015}{Causal Inference for Statistics, Social, and Biomedical Sciences}{\CUP.}
\obra{Greene}{2018}{Econometric Analysis}{8.ª ed., Pearson.} \recomendado
\obra{Heckman}{1979}{Sample Selection Bias as a Specification Error}{\textit{Econometrica}.}
\obra{Rosenbaum y Rubin}{1983}{The Central Role of the Propensity Score in Observational Studies for Causal Effects}{\textit{Biometrika}.}
\obra{Cinelli y Hazlett}{2020}{Making Sense of Sensitivity: Extending Omitted Variable Bias}{\textit{Journal of the Royal Statistical Society B}.}
\obra{Oster}{2019}{Unobservable Selection and Coefficient Stability}{\textit{Journal of Business \& Economic Statistics}.}
\end{bibliografia}

%% ---------------------------------------------------------------------
\modulo{3}{Econometría de panel avanzada}{2 h}
\begin{sesiones}
\sesion{7} Modelo de efectos aleatorios; evaluación y selección de modelos
  (Hausman, Mundlak).
\sesion{8} Error de medición y efectos fijos; paneles desbalanceados y
  atrición.
\end{sesiones}
\begin{bibliografia}
\obra{Wooldridge}{2010}{Econometric Analysis of Cross Section and Panel Data}{2.ª ed., MIT Press.}
\obra{Cameron y Trivedi}{2005}{Microeconometrics: Methods and Applications}{\CUP.}
\obra{Greene}{2018}{Econometric Analysis}{8.ª ed., Pearson.} \recomendado
\AngristMIT{Lecture Note 12.}
\end{bibliografia}

%% ---------------------------------------------------------------------
\modulo{4}{Diferencias en diferencias avanzado y control sintético}{4 h}
\begin{sesiones}
\sesion{9} DiD con covariables y con tratamientos continuos o de intensidad
  variable.
\sesion{10} Sensibilidad a pre-tendencias (Rambachan--Roth) y laboratorio
  comparativo de estimadores.
\sesion{11} Control sintético: pocas unidades tratadas, construcción del
  contrafactual y pruebas de placebo.
\sesion{12} DiD sintético: cuándo combina lo mejor de ambos enfoques.
\end{sesiones}
\begin{bibliografia}
\obra{Cunningham}{2021}{Causal Inference: The Mixtape}{Yale University Press.}
\obra{Rambachan y Roth}{2023}{A More Credible Approach to Parallel Trends}{\textit{Review of Economic Studies}.}
\obra{Borusyak, Jaravel y Spiess}{2024}{Revisiting Event-Study Designs: Robust and Efficient Estimation}{\textit{Review of Economic Studies}.}
\obra{Abadie, Diamond y Hainmueller}{2010}{Synthetic Control Methods for Comparative Case Studies}{\textit{Journal of the American Statistical Association}.}
\obra{Arkhangelsky, Athey, Hirshberg, Imbens y Wager}{2021}{Synthetic Difference-in-Differences}{\textit{American Economic Review}.}
\end{bibliografia}

%% ---------------------------------------------------------------------
\modulo{5}{Experimentos avanzados}{3 h}
\begin{sesiones}
\sesion{13} Poder y tamaño de muestra: laboratorio de simulación para diseñar
  un experimento.
\sesion{14} Experimentos con clusters y estratificación; inferencia por
  aleatorización.
\sesion{15} Heterogeneidad de efectos, múltiples pruebas y pre-registro.
\end{sesiones}
\begin{bibliografia}
\obra{List}{próx.}{Experimental Economics: Theory and Practice}{Aplicado a experimentos aleatorizados.}
\obra{List}{2022}{The Voltage Effect}{Currency.}
\obra{Abadie, Athey, Imbens y Wooldridge}{2023}{When Should You Adjust Standard Errors for Clustering?}{\textit{Quarterly Journal of Economics}.}
\obra{Angrist y Pischke}{2009}{Mostly Harmless Econometrics}{\PUP.}
\end{bibliografia}

%% ---------------------------------------------------------------------
\modulo{6}{Variables instrumentales y regresión discontinua avanzados}{3 h}
\begin{sesiones}
\sesion{16} IV avanzado: instrumentos débiles, diseños de jueces y
  \textit{shift-share}.
\sesion{17} RD difusa: ancho de banda, pruebas de manipulación y validez.
\sesion{18} Taller de RD con umbrales de scoring o regulatorios (datos
  sintéticos).
\end{sesiones}
\begin{bibliografia}
\obra{Angrist y Pischke}{2009}{Mostly Harmless Econometrics}{\PUP.}
\obra{Cunningham}{2021}{Causal Inference: The Mixtape}{Yale University Press.}
\obra{Imbens y Angrist}{1994}{Identification and Estimation of Local Average Treatment Effects}{\textit{Econometrica}.}
\obra{Imbens y Lemieux}{2008}{Regression Discontinuity Designs: A Guide to Practice}{\JoE.}
\end{bibliografia}

%% ---------------------------------------------------------------------
\modulo{7}{Aprendizaje automático causal}{2 h}
\begin{sesiones}
\sesion{19} \textit{Double/Debiased Machine Learning}: controles de alta
  dimensión y sobreajuste.
\sesion{20} Efectos heterogéneos con métodos de aprendizaje automático (por
  ejemplo, bosques causales).
\end{sesiones}
\begin{bibliografia}
\obra{Chernozhukov et al.}{2024}{Causal ML / Double Machine Learning}{\CUP.}
\end{bibliografia}

%% =====================================================================
\taller{3}{Análisis geoespacial y actividad de aplicación}
  {13 horas \textperiodcentered{} 13 sesiones de aproximadamente 1 hora \textperiodcentered{} requiere haber tomado los Talleres~1 y~2}

%% ---------------------------------------------------------------------
\modulo{1}{Análisis geoespacial}{7 h}
\begin{sesiones}
\sesion{1} Introducción conceptual: qué es y para qué sirve en el análisis
  económico (breve).
\sesion{2} Datos espaciales: puntos, líneas, polígonos y ráster.
\sesion{3} Sistemas de coordenadas y proyecciones.
\sesion{4} Distancia, proximidad y accesibilidad (p.~ej., sucursales y cajeros).
\sesion{5} Autocorrelación espacial.
\sesion{6} \textit{Hot spots} y análisis local.
\sesion{7} Modelos de regresión espacial.
\end{sesiones}
\begin{bibliografia}
\obra{Lovelace, Nowosad y Muenchow}{2025}{Geocomputation with R}{2.ª ed., \CRC.}
\obra{Pebesma y Bivand}{2023}{Spatial Data Science}{\CRC.}
\end{bibliografia}

%% ---------------------------------------------------------------------
\modulo{2}{Evaluación con dimensión espacial}{3 h}
\begin{sesiones}
\sesion{8} Efectos indirectos y exposición por distancia.
\sesion{9} Regresión discontinua geográfica y DiD con unidades espaciales.
\sesion{10} Caso práctico de acceso a servicios financieros con datos
  espaciales (sintéticos o públicos).
\end{sesiones}
\begin{bibliografia}
\obra{Pebesma y Bivand}{2023}{Spatial Data Science}{\CRC.}
\obra{Cunningham}{2021}{Causal Inference: The Mixtape}{Yale University Press.}
\end{bibliografia}

%% ---------------------------------------------------------------------
\modulo{3}{Actividad con aplicación a temas de la gerencia}{3 h}
\begin{sesiones}
\item[Sesiones 11 y 12] Actividad guiada: cada equipo aplica un diseño del
  Taller~1 o~2 a una pregunta de la gerencia, con uso de IA generativa en el
  flujo de trabajo.
\sesion{13} Presentaciones de resultados y retroalimentación.
\end{sesiones}
\begin{bibliografia}
\obra{Huntington-Klein}{2022}{The Effect}{\CRC.}
\obra{Cunningham}{2021}{Causal Inference: The Mixtape}{Yale University Press.}
\end{bibliografia}

\end{document}
